%% LyX 2.5.1 created this file.  For more info, see https://www.lyx.org/.
%% Do not edit unless you really know what you are doing.
\documentclass[11pt,brazil]{article}
\usepackage{lmodern}
\usepackage{lmodern}
\usepackage[T1]{fontenc}
\usepackage{textcomp}
\usepackage[utf8]{luainputenc}
\usepackage{color}
\definecolor{codebg}{rgb}{0.972656, 0.980469, 0.988281}
\definecolor{codeframe}{rgb}{0.886719, 0.910156, 0.941406}
\definecolor{commentgreen}{rgb}{0.0820312, 0.390625, 0.199219}
\definecolor{keywordblue}{rgb}{0.109375, 0.300781, 0.847656}
\definecolor{primary}{rgb}{0.113281, 0.246094, 0.6875}
\definecolor{secondary}{rgb}{0.0546875, 0.0859375, 0.160156}
\usepackage{colortbl}
\usepackage{babel}
\usepackage{cprotect}
\usepackage{booktabs}
\usepackage{amsmath}
\usepackage{amssymb}
\usepackage[a4paper]{geometry}
\geometry{verbose}
\usepackage[bookmarks=false,
 breaklinks=false,pdfborder={0 0 1},backref=false,colorlinks=true]
 {hyperref}
\hypersetup{pdftitle={Relatório Técnico - Biblioteca Modular em C},
 pdfauthor={Grupo de Desenvolvimento},
 linkcolor=primary,urlcolor=primary,citecolor=primary}

\makeatletter

%%%%%%%%%%%%%%%%%%%%%%%%%%%%%% LyX specific LaTeX commands.
%% Because html converters don't know tabularnewline
\providecommand{\tabularnewline}{\\}

%%%%%%%%%%%%%%%%%%%%%%%%%%%%%% User specified LaTeX commands.


% --- PACOTES DE LINGUAGEM E FONTE (Garante cópia perfeita de acentos no PDF) ---
\usepackage{cmap}
\PassOptionsToPackage{brazilian}{babel}
\usepackage{babel}


% --- PACOTES MATEMÁTICOS E ESTRUTURAIS ---
\usepackage{amsfonts}
\usepackage{xcolor}
\usepackage{listings}


% --- DEFINIÇÃO DE CORES ---







% --- CONFIGURAÇÃO DE EXIBIÇÃO DE CÓDIGO (LISTINGS) ---
\lstset{
    language=C,
    basicstyle=\ttfamily\small,
    backgroundcolor=,
    rulecolor=,
    keywordstyle=\bfseries,
    commentstyle=\itshape,
    frame=single,
    breaklines=true,
    captionpos=b,
    keepspaces=true,
    showstringspaces=false,
    tabsize=2,
    extendedchars=true
}



% --- DADOS DO DOCUMENTO ---
\title{\textbf{Relatório Técnico: Desenvolvimento da Biblioteca Modular \texttt{libmathgeo} em C}\\
\large Trabalho Prático 2 --- Estrutura, Compilação, Testes e Versionamento}
\author{\textbf{Programa de Graduação em Ciência da Computação}\\
Desenvolvimento de Software \& Programação Modular}
\date{\today}

\makeatother

\usepackage{listings}
\renewcommand{\lstlistingname}{Listagem}

\begin{document}
\title{\textbf{Relatório Técnico: Desenvolvimento de Biblioteca Modular Avançada
de Tópicos Matemáticos}}
\author{André De Vlieger Lima (16992758)\\
Lucas Correa Meira (15695588)\\
Vitor Wallace de Moraes Alves (8610353)}
\date{19 de Outubro de 2026}
\maketitle

\section{Descrição do Problema e do Domínio Escolhido}

Em sistemas computacionais de alto desempenho, engines gráficas e
simulações científicas, a modularidade e a reusabilidade de código
são requisitos indispensáveis. O problema abordado consiste na criação
de um ecossistema coeso de ferramentas matemáticas e geométricas na
linguagem C, estruturado na forma de uma biblioteca reutilizável seguindo
boas práticas de engenharia de software.

O domínio escolhido abrange cinco subáreas inter-relacionadas da matemática
aplicada e da computação gráfica:
\begin{enumerate}
\item \textbf{Álgebra Linear e Operações Vetoriais} (\texttt{vetores}):
manipulação de vetores estáticos e dinâmicos em 2D, 3D e$N$ dimensões;
\item \textbf{Geometria Analítica 2D/3D} (\texttt{geometria}): cálculo de
áreas, perímetros, centroides e volumes em formas regulares e malhas/polígonos
irregulares;
\item \textbf{Cálculo Diferencial e Integral} (\texttt{calculo}): diferenciação
finita, vetores gradientes e integração numérica simples e dupla;
\item \textbf{Métodos Numéricos} (\texttt{calculo\_numerico}): localização
de raízes de equações não lineares através de métodos iterativos com
tratamento formal de erros;
\item \textbf{Polinômios Esparsos} (\texttt{polinomios}): representação
e manipulação eficiente de polinômios de alto grau utilizando listas
encadeadas dinâmicas.
\end{enumerate}

\section{Arquitetura do Software}

A arquitetura do projeto adota o princípio do encapsulamento estrito
e da visibilidade mínima. A estrutura física de diretórios separa
os contratos de interface pública, o código de implementação, os testes
automatizados e a aplicação cliente final:

\bgroup
\begin{lstlisting}[language=bash]
projeto/
|-- include/           # Interface publica (.h)
|   |-- vetores.h
|   |-- geometria.h
|   |-- calculo.h
|   |-- calculo_numerico.h
|   `-- polinomios.h
|-- src/               # Implementacao interna (.c)
|   |-- vetores.c
|   |-- geometria.c
|   |-- calculo.c
|   |-- calculo_numerico.c
|   `-- polinomios.c
|-- tests/             # Testes automatizados
|   `-- test_main.c
|-- app/               # Aplicacao cliente
|   `-- main.c
|-- build/             # Objetos intermediarios (.o)
|-- lib/               # Bibliotecas geradas (.a, .so)
|-- bin/               # Executaveis finais
|-- Makefile           # Automacao de compilacao
`-- README.md          # Manual de uso
\end{lstlisting}
\leavevmode\egroup{}

\section{Organização e Responsabilidade dos Módulos}

A biblioteca é composta por 5 módulos funcionalmente distintos:

\begin{table}[h!]
\centering\caption{Módulos da biblioteca e suas responsabilidades}
\label{tab:modulos}%
\begin{tabular}{lll}
\toprule 
\textbf{Módulo} & \textbf{Arquivos} & \textbf{Responsabilidade Funcional}\tabularnewline
\midrule 
\texttt{vetores} & \texttt{vetores.h / .c} & Operações vetoriais, normas, produtos escalar/vetorial em 2D/3D/ND.\tabularnewline
\texttt{geometria} & \texttt{geometria.h / .c} & Geometria computacional: Shoelace, Ray-Casting e volumes 3D.\tabularnewline
\texttt{calculo} & \texttt{calculo.h / .c} & Diferenciação centrada, gradientes, trapézios e ponto médio 2D.\tabularnewline
\texttt{calculo\_numerico} & \texttt{calculo\_numerico.h / .c} & Busca de raízes via Bissecção e Newton-Raphson com enums de erro.\tabularnewline
\texttt{polinomios} & \texttt{polinomios.h / .c} & Polinômios esparsos com alocação dinâmica e listas encadeadas.\tabularnewline
\bottomrule
\end{tabular}
\end{table}


\section{Descrição da API Pública}

A API pública exporta mais de 20 funções através dos arquivos de cabeçalho.
Os detalhes de implementação permanecem restritos aos arquivos de
código-fonte (\texttt{.c}).

\subsection{Módulo\texttt{vetores.h}}
\begin{itemize}
\item \texttt{VetorND vec\_criar\_nd(int dimension, const double {*} components)}:
aloca vetor$N$-dimensional dinâmico;
\item \texttt{void vec\_liberar\_nd(VetorND {*} v)}: desaloca a memória
do vetor dinâmico;
\item \texttt{Vetor3D vec\_somar3d(Vetor3D v1, Vetor3D v2)}: soma de vetores
em$\mathbb{R}^{3}$;
\item \texttt{double vec\_produto\_escalar3d(Vetor3D v1, Vetor3D v2)} e\texttt{vec\_produto\_vetorial3d}:
produto escalar e vetorial.
\end{itemize}

\subsection{Módulo\texttt{geometria.h}}
\begin{itemize}
\item \texttt{double geo\_area\_poligono\_irregular(PoligonoIrregular2D
p)}: área por Fórmula de Shoelace;
\item \texttt{int geo\_ponto\_em\_poligono\_irregular(Ponto2D ponto, PoligonoIrregular2D
p)}: algoritmo de Ray-Casting;
\item \texttt{double geo\_volume\_malha\_fechada(MalhaIrregular3D m)}: volume
via decomposição tetraédrica.
\end{itemize}

\subsection{Módulos\texttt{calculo.h} e\texttt{calculo\_numerico.h}}
\begin{itemize}
\item \texttt{double calc\_derivada\_1d(Funcao1D f, double x, double h)}:
aproximação por diferenças finitas centradas;
\item \texttt{void calc\_gradiente(FuncaoND f, const double {*} ponto, int
num\_vars, double h, double {*} res)}: vetor gradiente;
\item \texttt{StatusCalculo num\_bisseccao(...)} e\texttt{num\_newton\_raphson(...)}:
solução iterativa de raízes com retorno de\texttt{StatusCalculo}.
\end{itemize}

\subsection{Módulo\texttt{polinomios.h}}
\begin{itemize}
\item \texttt{Polinomio poli\_criar\_vazio(void)}: inicializa polinômio
nulo;
\item \texttt{void poli\_adicionar\_termo(Polinomio {*} p, double coef,
int exp)}: insere termo ordenado por expoente;
\item \texttt{Polinomio poli\_derivada\_algebrica(Polinomio p)}: gera a
derivada analítica exata.
\end{itemize}

\section{Descrição da Aplicação Cliente}

A aplicação cliente, implementada no arquivo app/main.c, tem como
objetivo demonstrar a utilização prática e integrada das funcionalidades
fornecidas pela biblioteca modular. Ela funciona como um programa
executável independente que consome a API pública projetada.

\subsection{Respeito Estrito à API Pública}

Conforme exigido pelos requisitos de arquitetura e encapsulamento
do projeto, a aplicação cliente:
\begin{itemize}
\item Apenas os cabeçalhos presentes no diretório {[} include/ {]} (vetores.h,
polinomios.h, geometria.h, calculo.h e calculo\_numerico.h) são importados
via \#include;
\item Em nenhuma hipótese o arquivo app/main.c faz inclusão direta de arquivos
localizados no diretório {[} src/ {]};
\item Não acessa estruturas internas/privadas: toda a manipulação de dados
ocorre por meio das funções públicas expostas e dos tipos definidos
na interface pública.
\end{itemize}

\section{Processo de Compilação Utilizando Makefile}

O build do projeto é gerido pelo\texttt{Makefile}, contendo os seguintes
alvos padronizados:
\begin{itemize}
\item \texttt{make all}: Compila os módulos, gera as bibliotecas estática/dinâmica
e o executável cliente;
\item \texttt{make clean}: Remove todos os artefatos gerados nas pastas\texttt{build/},\texttt{lib/}
e\texttt{bin/};
\item \texttt{make run}: Compila se necessário e executa a aplicação;
\item \texttt{make test}: Compila e executa a suíte de testes automatizados;
\item \texttt{make install PREFIX=\$HOME/.local}: Instala cabeçalhos e bibliotecas
no diretório de destino.
\end{itemize}

\section{Estratégia de Compilação Incremental}

As dependências são mapeadas de forma que cada objeto em\texttt{build/\%.o}
depende exclusivamente de seu arquivo fonte\texttt{src/\%.c} e do
respectivo cabeçalho em\texttt{include/\%.h}. Alterar um módulo recompila
apenas o arquivo alterado e refaz o link, poupando tempo de compilação.

\section{Geração e Utilização da Biblioteca Estática}

A biblioteca estática (\texttt{libapp.a}) é criada via utilitário\texttt{ar}:

\bgroup
\begin{lstlisting}[language=bash]
gcc -Wall -Wextra -Iinclude -c src/vetores.c -o build/vetores.o
gcc -Wall -Wextra -Iinclude -c src/polinomios.c -o build/polinomios.o
gcc -Wall -Wextra -Iinclude -c src/geometria.c -o build/geometria.o
gcc -Wall -Wextra -Iinclude -c src/calculo.c -o build/calculo.o
gcc -Wall -Wextra -Iinclude -c src/calculo_numerico.c -o build/calculo_numerico.o
ar rcs lib/libapp.a build/*.o
gcc -Iinclude app/main.c lib/libapp.a -lm -o bin/app_estatico

./bin/app_estatico
\end{lstlisting}
\leavevmode\egroup{}

\section{Geração e Utilização da Biblioteca Dinâmica}

A biblioteca dinâmica (\texttt{libapp.so}) exige a opção\texttt{-fPIC}
e a flag\texttt{-shared}:

\bgroup
\begin{lstlisting}[language=bash]
gcc -Wall -Wextra -fPIC -Iinclude -c src/*.c
# Move os .o gerados para a pasta build (opcional, dependendo de como organiza)
mv *.o build/
gcc -shared build/*.o -o lib/libapp.so
gcc -Iinclude app/main.c -Llib -lapp -lm -o bin/app_dinamico
LD_LIBRARY_PATH=lib ./bin/app_dinamico
\end{lstlisting}
\leavevmode\egroup{}

\section{Comparação entre Ligação Estática e Dinâmica}

Este projeto implementa e analisa as duas abordagens fundamentais:
ligação estática (libapp.a) e ligação dinâmica (libapp.so).

\subsection{Biblioteca Estática (libapp.a)}

A biblioteca estática é uma coleção de arquivos objeto (.o) empacotados
em um único arquivo de arquivo, utilizando o utilitário\emph{ar}.
\begin{itemize}
\item Mecanismo: durante o processo de compilação, o linker extrai os símbolos
necessários da biblioteca e os copia diretamente para dentro do binário
executável final;
\item Independência em Runtime: o executável resultante é totalmente autônomo.
Ele pode ser movido para outro ambiente Linux idêntico e executado
sem a necessidade de manter o arquivo libapp.a no sistema.
\end{itemize}

\subsection{Biblioteca Dinâmica (libapp.so)}

A biblioteca dinâmica (Shared Object) é gerada compilando o código-fonte
com a flag de código independente de posição (-fPIC) e vinculada via
gcc -shared.
\begin{itemize}
\item Mecanismo: O linker não copia o código da biblioteca para o executável;
em vez disso, insere tabelas de referência a símbolos externos. A
resolução real dos endereços das funções é delegada ao carregador
do sistema operacional (dynamic linker) no momento em que o programa
é iniciado;
\item Dependência em Runtime: O executável necessita que o arquivo .so esteja
presente no disco e localizado em diretórios conhecidos pelo sistema
ou especificados pela variável de ambiente LD\_LIBRARY\_PATH.
\end{itemize}

\subsection{Resultados e Demonstração Prática no Projeto}

Durante o desenvolvimento do trabalho, as seguintes evidências experimentais
comprovaram as diferenças teóricas:

\subsubsection{Inspeção de Dependências com ldd:}

Ao executar o comando ldd bin/client na versão com ligação dinâmica,
o sistema listou explicitamente a dependência de tempo de execução:
\begin{enumerate}
\item libapp.so => lib/libapp.so (0x00007f...)
\item libm.so.6 => /lib/x86\_64-linux-gnu/libm.so.6 (0x00007f...)
\item libc.so.6 => /lib/x86\_64-linux-gnu/libc.so.6 (0x00007f...)
\end{enumerate}

\subsubsection{O “Teste da Remoção” de Diretório:}

Executável Dinâmico: Ao remover ou mover temporariamente o diretório
lib e tentar executar ./bin/client, o sistema operacional interrompeu
a execução imediatamente exibindo o erro: error while loading shared
libraries: libapp.so: cannot open shared object file.

Executável Estático: Ao compilar um binário estático via: gcc -Iinclude
app/main.c lib/libapp.a -lm -o bin/client\_estatico e em seguida remover
a pasta lib/, o programa executou perfeitamente sem falhar, comprovando
que todo o código da biblioteca já estava incorporado ao executável.

\section{Estratégia e Resultados dos Testes}

\subsection{Estratégia de Teste}

A estratégia adotada apoia-se em três pilares principais:
\begin{enumerate}
\item Isolamento Modular: cada módulo da biblioteca possui um arquivo de
teste dedicado (test\_vetores.c, test\_polinomios.c, test\_geometria.c,
test\_calculo.c e test\_calculo\_num. O orquestrador central test\_main.c
executa sequencialmente as suítes de cada módulo;
\item Validação por Asserções (assert.h): utilizam-se asserções em C para
verificar condições lógicas de contorno, estados de erro e propriedades
matemáticas esperadas;
\item Comparação com Tolerância Numérica: devido às imprecisões inerentes
à repre- sentação de ponto flutuante em computadores, operações numéricas
não utilizam igualdade estrita (==). Em vez disso, implementou-se
a verificação com uma tolerância =$10^{-4}$ via função fabs().
\end{enumerate}

\subsection{Execução Automatizada via Makefile}

A automação é integrada ao processo de construção do projeto através
do alvo make test. O comando compila os arquivos de teste vinculando-os
à biblioteca estática e à biblioteca matemática (-Llib -lapp -lm)
e executa o binário resultante bin/test.

\section{Exemplo de Teste que Inicialmente Falhou e Posterior Correção}

Em conformidade com as boas práticas de desenvolvimento orientado
a testes (TDD) e os requisitos do edital do projeto, registra-se nesta
seção um caso real de falha identificada durante a execução da suíte
de testes e o respectivo processo de correção.

\subsection{Descrição do Erro Encontrado}

Durante a execução inicial da suíte de testes do Módulo de Cálculo
Diferencial e Integral, a asserção responsável por validar o vetor
gradiente da função f (x, y) = x2 + y2 no ponto (1.0, 1.0) falhou.
A derivada parcial esperada em relação a x no ponto x = 1.0 é fx(1,
1) = 2 · (1.0) = 2.0.

No entanto, o algoritmo retornou o valor \textminus 2.0, ocasionando
a interrupção da suíte pelo mecanismo de asserção.

\subsection{Saída do Terminal Acusando a Falha}

Ao executar o comando make test, a suíte apresentou a seguinte mensagem
de erro no terminal:

\bgroup
\begin{lstlisting}
--- Executando Suíte de Testes ---
[MÓDULO] Vetores...
Aprovado: Vetores 3D e N-D
[MÓDULO] Polinômios Esparsos...
Aprovado: Polinômios
[MÓDULO] Geometria Analítica 2D/3D...
Aprovado: Polígonos 2D e Geometria Espacial
[MÓDULO] Cálculo Diferencial e Integral...
test: tests/test_calculo.c:26: testar_calculo: Assertion `fabs((grad[0]) - (2.0)) < (1e-
4)' failed.
make: *** [Makefile:41: test] Abortado (arquivo core criado)
\end{lstlisting}
\leavevmode\egroup{}

\subsection{Causa Raiz da Falha}

Ao inspecionar o arquivo src/calculo.c, identificou-se um erro de
sinal na implementação da aproximação por diferenças finitas centradas
na função calc\_derivada\_parcial:

\[
f'(x)\approx\frac{f(x+h)-f(x-h)}{2h}
\]

No código original, a ordem dos termos no numerador estava invertida,
calculando$f(x-h)-f(x+h)$, o que resultava no sinal oposto da derivada:

\bgroup
\begin{lstlisting}
/* Código INCORRETO em src/calculo.c */
p[var_index] = ponto[var_index] + h;
double l_superior = f(p_sup, num_vars);

p[var_index] = ponto[var_index] - h;
double l_inferior = f(p_inf, num_vars);

// return (f(p,num_vars) - f(p,num_vars))/(2*h); por conta da alocação dinamica e do free fui obrigado copiar os valores pra f_mais e f_menos
    free(p);
return (l_inferior - l_superior)/(2.0 * h); // ERRO: Ordem invertida!
\end{lstlisting}
\leavevmode\egroup{}

\subsection{Correção Aplicada}

A função foi corrigida no arquivo src/calculo.c, ajustando-se a subtração
para a ordem correta$f(x-h)-f(x+h)$:

\bgroup
\begin{lstlisting}
/* Código CORRIGIDO em src/calculo.c */
p[var_index] = ponto[var_index] + h;
double l_superior = f(p_sup, num_vars);

p[var_index] = ponto[var_index] - h;
double l_inferior = f(p_inf, num_vars);

// return (f(p,num_vars) - f(p,num_vars))/(2*h); por conta da alocação dinamica e do free fui obrigado copiar os valores pra f_mais e f_menos
    free(p);
return (l_superior - l_inferior)/(2.0 * h); // CORRETO
\end{lstlisting}
\leavevmode\egroup{}

\subsection{Validação Pós-Correção}

Após recompiar a biblioteca e reexecutar a suíte de testes com make
clean \&\& make test, todos os testes passaram com 100\% de sucesso:

\bgroup
\begin{lstlisting}
--- Executando Suíte de Testes ---
[MÓDULO] Vetores...
Aprovado: Vetores 3D e N-D
[MÓDULO] Polinômios Esparsos...
Aprovado: Polinômios
[MÓDULO] Geometria Analítica 2D/3D...
Aprovado: Polígonos 2D e Geometria Espacial
[MÓDULO] Cálculo Diferencial e Integral...
Aprovado: Derivadas, Gradientes e Integrais 1D/2D
[MÓDULO] Cálculo Numérico...
Aprovado: Bissecção e Newton-Raphson

Todos os testes passaram com sucesso!
\end{lstlisting}
\leavevmode\egroup{}

\section{Estratégia de Versionamento Utilizando Git}

A estratégia de versionamento adotada visou manter a integridade da
base de código, garantir a rastreabilidade das alterações e documentar
a evolução modular da biblioteca, seguindo sempre as melhores práticas
de\emph{Atomicidade dos Commits},\emph{Mensagens de Commit Claras}
e\emph{Uso do .gitignore}.

\section{Descrição de Branch, Merge e Tag v1.0}

Decidiu-se seguir da proposta: cada módulo teria de ser uma Branch
separada; dessa forma, poderia-se facilmente delimitar as tarefas
do projeto e analisar a questão de falhas no código, já que cada branch
teve de se encapsular num pull request que, em fim, lida a um merge.

Seguindo as boas práticas, depois do desenvolvimento do projeto, colocou
uma tag v1.0 para o github do mesmo, indicando que o desenvolvimento
havia se encerado.

\section{Dificuldades Encontradas e Decisões de Projeto}

O desenvolvimento da biblioteca modular exigiu decisões cuidadosas
de arquitetura para harmonizar desempenho, clareza de código e encapsulamento
em linguagem C.

\subsection{Dificuldades e Soluções}
\begin{enumerate}
\item Resolução de Caminho da Biblioteca Dinâmica em Runtime (LD\_LIBRARY\_PATH):
\begin{enumerate}
\item Problema: ao executar o cliente compilado com ligação dinâmica (./bin/client),
o carregador do sistema operacional falhava ao localizar a biblioteca
libapp.so;
\item Decisão de Projeto: para fins de testes e execução rápida, utilizou-se
a variável de ambiente export LD\_LIBRARY\_PATH=./lib:\$LD\_LIBRARY\_PATH.
Para o alvo automatizado make run, optou-se por embutir o caminho
de busca em tempo de execução via flag do linker (-Wl,-rpath=./lib),
permitindo a execução imediata do programa pelo usuário.
\end{enumerate}
\item Gerenciamento de Memória em Polinômios Esparsos e Vetores N-D:
\begin{enumerate}
\item Problema: risco de vazamento de memória (memory leak) ao manipular
listas encadeadas dinâmicas e vetores N -dimensionais alocados com
malloc;
\item Decisão de Projeto: estabeleceu-se o princípio da responsabilidade
explícita na API pública, onde cada função construtora (vec\_criar\_nd,
poli\_criar\_vazio) obrigatoriamente possui uma função correspondente
de destruição/liberação (vec\_liberar\_nd, poli\_liberar), devidamente
testadas com verificações de ponteiros nulos.
\end{enumerate}
\item Comparação de Valores de Ponto Flutuante nos Testes
\begin{enumerate}
\item Problema: comparações diretas de igualdade (==) entre valores do tipo
double falhavam devido ao arredondamento;
\item Decisão de Projeto: foi adotado a padronização de testes numéricos
utilizando uma margem de tolerância =$10^{-4}$ via macro ASSERT\_NEAR,
assegurando a robustez das asserções.
\end{enumerate}
\item Inversão de Sinal na Derivada Numérica:
\begin{enumerate}
\item Problema: conforme detalhado na Seção 12, a ordem dos operandos no
cálculo da diferença finita estava invertida, resultando em gradientes
com sinal trocado;
\item Decisão de Projeto: a correção foi guiada pela suíte de testes unitários
automatizados, comprovando a eficácia do uso de testes durante o ciclo
de desenvolvimento.
\end{enumerate}
\end{enumerate}

\section{Resultados Obtidos e Possíveis Melhorias}

\subsection{Resultados Obtidos}

O projeto atingiu plenamente todos os requisitos funcionais, não funcionais
e de arquitetura preestabelecida:
\begin{itemize}
\item Arquitetura Modular Completa: desenvolvimento bem-sucedido dos 5 módulos
matemáticos (vetores, polinomios, geometria, calculo e calculo\_numerico)
com estrito encapsulamento e separação entre interface (.h) e implementação
(.c);
\item Geração Dupla de Bibliotecas: automação completa no Makefile para
compilação e empacotamento da biblioteca estática (libapp.a) e biblioteca
dinâmica (libapp.so);
\item Aplicação Cliente Funcional: implementação da aplicação app/main.c
que consome exclusivamente a API pública dos módulos e valida o funcionamento
integrado da biblioteca;
\item Testes Automatizada com 100\% de Aprovação: verificação rigorosa de
todos os métodos numéricos e geométricos via make test;
\item Versionamento Rigoroso: histórico organizado com uso de branches,
merges explícitos e identificação da versão final via tag v1.0.
\end{itemize}

\subsection{Possíveis Melhorias Futuras}

Como perspectivas de evolução para versões futuras do projeto, destacam-se:
\begin{enumerate}
\item Paralelização de Cálculos Intensivos: utilização da biblioteca OpenMP
ou POSIX Threads para paralelizar o cálculo de integrais duplas e
gradientes N-dimensionais, reduzindo o tempo de processamento em CPUs
multicore;
\item Implementação de Métodos Numéricos de Ordem Superior: expansão do
módulo de cálculo numérico com métodos de maior convergência;
\item Geração Automática de Documentação: adoção do padrão de comentários
Doxygen nos arquivos de cabeçalho (.h) para possibilitar a geração
automática de manuais de referência em HTML e PDF;
\item Suíte de Benchmark de Desempenho: adição de um módulo de medição de
tempo de execução e perfilamento de memória (profiling) para comparar
a eficiência do código em diferentes níveis de otimização do compilador
(-O2, -O3).
\end{enumerate}

\end{document}
